\section{把评估拆开看：定义、直觉、机制、应用、局限}

前面已经把课程里出现的 benchmark 和问题铺开了。这里把它们重新整理成一个更稳定的分析框架，方便你以后面对一个新榜单时快速判断它到底在测什么。

\subsection{定义：评估到底是什么}

严格地说，评估是在一个给定问题设定下，把模型、系统或方法映射成可比较的数值或排序。这个映射必须先规定好四件事：
\begin{itemize}
    \item 任务输入是什么。
    \item 模型如何被调用。
    \item 输出怎么被度量。
    \item 分数怎么被解释。
\end{itemize}

\begin{figure}[H]
\centering
\begin{tikzpicture}[node distance=0.95cm, every node/.style={font=\small, align=center}]
\node[draw, rounded corners, fill=blue!10, minimum width=3.0cm, minimum height=0.8cm] (a) {goal};
\node[draw, rounded corners, fill=blue!18, right=1.2cm of a, minimum width=3.0cm, minimum height=0.8cm] (b) {inputs};
\node[draw, rounded corners, fill=blue!26, right=1.2cm of b, minimum width=3.0cm, minimum height=0.8cm] (c) {calling convention};
\node[draw, rounded corners, fill=blue!34, right=1.2cm of c, minimum width=3.0cm, minimum height=0.8cm] (d) {metric};
\draw[-{Latex[length=3mm]}, thick] (a) -- (b);
\draw[-{Latex[length=3mm]}, thick] (b) -- (c);
\draw[-{Latex[length=3mm]}, thick] (c) -- (d);
\end{tikzpicture}
\caption{评估不是一个分数，而是一条从目标到 metric 的定义链}
\end{figure}

\begin{knowledgebox}{定义层的判断标准}
如果一套评估没有把目标、输入、调用和 metric 分开写清楚，那么它通常不适合直接拿来比较模型。它可能有参考价值，但不该被当成最终答案。
\end{knowledgebox}

\subsection{直觉：为什么评估看起来简单，做起来很难}

直觉上，评估像是“把模型扔到题目里跑一遍”。实际问题是，语言模型不是一个固定函数，而是一个会被 prompt、工具、解码策略和上下文长度改变行为的系统。

\begin{table}[H]
\centering
\begin{tabular}{p{3.2cm}p{4.4cm}p{5.2cm}}
\toprule
\textbf{看起来像} & \textbf{实际更像} & \textbf{后果} \\
\midrule
算分 & 系统设计 & 不同 prompt 会改变排名 \\
测能力 & 选定场景下的能力投影 & benchmark 不一定代表真实使用 \\
看总分 & 看任务分解 & 单个分数会掩盖失败模式 \\
\bottomrule
\end{tabular}
\caption{评估的直觉与现实之间的落差}
\end{table}

\begin{importantbox}{为什么“看起来简单”是危险的}
评估一旦被简化成“跑脚本”，就容易忽略污染、偏差、提示词敏感性和 judge bias。越是看起来规整的数字，越需要怀疑它背后的定义是否足够稳。
\end{importantbox}

\subsection{机制：一个完整评估管线怎么运作}

可以把现代评估管线理解成五步：
\begin{enumerate}
    \item 选择任务集合。
    \item 设计输入格式和 prompt。
    \item 运行模型或系统。
    \item 收集输出并计算指标。
    \item 聚合结果并解释。
\end{enumerate}

\begin{figure}[H]
\centering
\begin{tikzpicture}[node distance=0.9cm, every node/.style={font=\small, align=center}]
\node[draw, rounded corners, fill=green!10, minimum width=3.3cm, minimum height=0.8cm] (t1) {task set};
\node[draw, rounded corners, fill=green!18, below=of t1, minimum width=3.3cm, minimum height=0.8cm] (t2) {prompt / input};
\node[draw, rounded corners, fill=green!26, below=of t2, minimum width=3.3cm, minimum height=0.8cm] (t3) {model or system};
\node[draw, rounded corners, fill=green!34, below=of t3, minimum width=3.3cm, minimum height=0.8cm] (t4) {metric computation};
\node[draw, rounded corners, fill=green!42, below=of t4, minimum width=3.3cm, minimum height=0.8cm] (t5) {interpretation};
\draw[-{Latex[length=3mm]}, thick] (t1) -- (t2);
\draw[-{Latex[length=3mm]}, thick] (t2) -- (t3);
\draw[-{Latex[length=3mm]}, thick] (t3) -- (t4);
\draw[-{Latex[length=3mm]}, thick] (t4) -- (t5);
\end{tikzpicture}
\caption{评估管线的五个环节：任务、输入、系统、指标、解释}
\end{figure}

\begin{warningbox}{机制层最常见的误解}
很多争论其实不是“哪个模型更强”，而是“哪个评估管线更公平”。如果任务、prompt 或 judge 改了，结果就可能改，甚至改得比模型差异还大。
\end{warningbox}

\subsection{应用：什么时候该用什么评估}

不同场景应该用不同评估。

\begin{table}[H]
\centering
\begin{tabular}{p{3.0cm}p{4.0cm}p{5.8cm}}
\toprule
\textbf{场景} & \textbf{应该更关心什么} & \textbf{更合适的评估} \\
\midrule
买模型 / 选产品 & 用户体验、价格、稳定性 & arena、成本前沿、真实流量 \\
做基础研究 & 可解释的能力增长 & benchmark + ablation + scaling law \\
做安全审查 & 有害行为与越狱 & curated safety sets, red teaming \\
做系统优化 & 吞吐、延迟、成本 & throughput / latency / cost profiling \\
\bottomrule
\end{tabular}
\caption{评估方法必须随目标变化，而不是一套分数走天下}
\end{table}

\begin{knowledgebox}{如何快速选 benchmark}
先问自己：我是在测知识、指令遵循、工具使用、推理，还是安全风险？如果问题没想清楚，先别急着选分数，先选任务定义。
\end{knowledgebox}

\subsection{局限：为什么没有一个 benchmark 能包打天下}

任何 benchmark 都至少有四类局限：
\begin{itemize}
    \item \textbf{覆盖不全}：真实场景总比 benchmark 更杂。
    \item \textbf{可被游戏化}：一旦分数成为目标，就会有人围着分数优化。
    \item \textbf{会过时}：数据集一旦饱和，区分力就下降。
    \item \textbf{会偏置}：人类判断、LLM judge、流量数据都可能带偏。
\end{itemize}

\begin{importantbox}{局限不是 bug，而是事实}
评估本身就是近似。你不是在寻找一个绝对真理，而是在寻找一个足够稳、足够相关、足够可解释的代理指标。
\end{importantbox}

